\documentclass[11pt]{article}

\usepackage[margin=25mm]{geometry}

\usepackage{amsmath,amssymb,fontspec,unicode-math}

\IfFontExistsTF{DejaVu Serif}{\setmainfont{DejaVu Serif}}{\setmainfont{Latin Modern Roman}}

\IfFontExistsTF{Latin Modern Math}{\setmathfont{Latin Modern Math}}{\setmathfont{DejaVu Math TeX Gyre}}

\usepackage{xurl,hyperref}

\setlength{\parindent}{0pt}\setlength{\parskip}{6pt}

\begin{document}

\title{Fair Cooperation in Mixed-Motive Games via Conflict-Aware Gradient Adjustment}

\author{ELVIA preliminary AI analysis}\date{21 October 2026}\maketitle

Preliminary assessment: Major.\par

Independent scientific review remains pending. Author adoption has not been requested.

Bounded smooth counterexamples refute finite-step protected improvement, universal fairness convergence and the simultaneous-update extension. The reported empirical gains are not themselves disproved.

\section*{Review scope}

Rapid current-version reassessment of every substantive concern in the supplied LaTeX report: the implicated paper passages and relevant author-released source paths were inspected, existing proposed corrections were briefly checked, and impact was calibrated to the demonstrated scope. No author-provided or historical verification program was executed.

\section*{Problem 1: Orthogonal projection does not preserve an objective after a positive finite step}

Source: Theorem 3.1, PDF p. 5; Theorem B.2, PDF p. 14.

Assessment: supported; impact: major

\subsection*{Source claim (paraphrase)}

Paraphrase: Theorem 3.1, p. 5, and its restatement as Theorem B.2, p. 14, assert that differentiable \(L\)-smooth individual and collective objectives both do not decrease under the FCGrad update for sufficiently small \(\eta>0\). The conflict-case proof uses \(b^\top g=0\) to infer preservation of the collective objective.

\subsection*{Analysis}

At \((0,0)\), take bounded \(1\)-smooth objectives \(I(x,y)=\sin x+\sin y\) and \(C(x,y)=1-\sin x+\cos y\). Their gradients are \((1,1)\) and \((-1,0)\); conflict projection leaves direction \((0,1)\), orthogonal to \(\nabla C\). Nevertheless \(C(0,\eta)=1+\cos\eta<2=C(0,0)\) for every sufficiently small \(\eta>0\). A zero directional derivative does not ensure a nonnegative finite-step change.

\subsection*{Finding and scope}

The theorem's finite-step protection is false under its smoothness assumptions. The example leaves observed empirical gains unadjudicated.

\subsection*{Proposed correction}

The existing direction supports first-order protection only: smoothness gives \(C(\theta+\eta g)-C(\theta)\geq-L\eta^2\|g\|^2/2\). A separate proposed common-ascent modification applies only to noncollinear conflicting gradients: for \(q=a^\top b<0\), let \(h=a-qb/\|b\|^2\), \(m=\|h\|^2>0\), and \(d=h+mb/[2(-q)]\). Then \(a^\top d=m/2>0\), \(b^\top d>0\); a sufficiently small step improves both fixed-agent objectives. This changes the algorithm and does not establish fairness convergence or simultaneous-agent improvement.

A local correction or qualification is proposed where specified; author adoption is not established.

\section*{Problem 2: Recurrent conflicts can leave a nonzero value gap forever}

Source: Theorem 3.2, PDF p. 5; Theorem B.5, PDF p. 16.

Assessment: supported; impact: major

\subsection*{Source claim (paraphrase)}

Paraphrase: Theorem 3.2, p. 5, and Theorem B.5, p. 16, assume \(L\)-smooth values, \(\delta_t=V_{\rm ind}(\theta_t)-V_{\rm col}(\theta_t)\), positive steps satisfying

\[\eta_t\leq|\delta_t|/L,\qquad \sum_t\eta_t=\infty,\qquad \sum_t\eta_t^2<\infty,\]

\subsection*{Analysis}

Let \(I(\theta)=\sin\theta\), \(C(\theta)=1-\sin\theta\), and \(\theta_0=0\). Their opposing gradients \(1,-1\) project to zero, so every iterate remains zero and the value gap stays \(-1\). Steps \(\eta_t=1/(t+1)\) satisfy the upper bound, divergent sum and square-summable condition with \(L=1\). Conflict recurs forever without progress toward equal values.

\subsection*{Finding and scope}

The universal fairness-convergence conclusion requires a nondegeneracy or progress assumption absent from the stated theorem.

\section*{Problem 3: Local guarantees do not imply individual improvement under simultaneous multi-agent updates}

Source: Paragraph following Theorem 3.1 and Section 3.1 single-agent parameters, PDF p. 5.

Assessment: supported; impact: major

\subsection*{Source claim (paraphrase)}

The paragraph immediately after Theorem 3.1, p. 5, extends the single-agent result to all agents when all use FCGrad. The mathematical objects in Section 3.1 are gradients with respect to the parameters of one agent.

\subsection*{Analysis}

Take returns \(R_1=\sin\theta_1-3\sin\theta_2\) and \(R_2=3\sin\theta_1+4\sin\theta_2\), with collective return \((R_1+R_2)/2\). At zero, with the allowed mixing weight \(\beta=1/2\), the nonconflict FCGrad directions are \(1.5\) and \(2.25\). Each own-coordinate update improves its local objectives, but their simultaneous effect on \(R_1\) has derivative \(1.5-3(2.25)=-5.25\). The single-agent argument leaves these cross-agent terms uncontrolled.

\subsection*{Finding and scope}

The stated simultaneous-update extension does not follow and fails on this bounded smooth game. A separate multi-agent guarantee would need additional assumptions.

\section*{Limits}

No author-provided program, training procedure, model inference or wet-lab experiment was executed in this review.\par

Historical execution results are retained only with that attribution; a source-level defect alone does not identify the program or configuration used for the published experiments.\par

No author contact, adoption, independent human scientific review or novelty determination is established.\par

The findings apply only to the identified claims, source versions and stated assumptions; they do not certify or invalidate the entire paper.\par

\section*{Source identity}

NeurIPS 2025 official proceedings PDF, freshly retrieved from the current publisher landing page

\url{https://proceedings.neurips.cc/paper_files/paper/2025/file/4830f08ffb59c0d49aef99f3a6f4f85f-Paper-Conference.pdf}

Reviewed PDF SHA-256: \nolinkurl{1ce755f50690904a2aa47cfeb72ef6b310ed85744b9e860ab5e8256adc27f9bb}

The exact inspected current paper PDF is identified by its SHA-256. The current publisher PDF was freshly downloaded and is byte-identical to the previously retained proceedings copy. Latest author-manuscript history checked: 2508.17696v1 (last updated 2025-08-25T06:08:13Z). The current formal proceedings remain the primary manuscript; the latest indexed preprint does not postdate or replace the formal work in this channel. Current author resource: https://github.com/wjkim1202/fcgrad at commit 50f778c7d8c0c22323c38070d7019fd8826c2dfb.  An exact PDF acquisition timestamp is not established in this reassessment; source-currentness checks have their own recorded timestamps and file modification times were not substituted.

\end{document}
